
{\Large\bfseries\color{accent}4. Multipoles, parity and selection rules}
\hrule\smallskip
\begin{multicols}{2}
\subsection*{Operators}
Electric multipole \(E k\): rank \(k\), parity \((-1)^k\).
Magnetic multipole \(M k\): rank \(k\), parity \((-1)^{k+1}\).
Orbital state parity is \((-1)^l\); multi-electron
configuration parity is \((-1)^{\sum_i l_i}\).
\subsection*{E1 / M1 / E2 at a glance}
\begin{tabular}{@{}lll@{}}\toprule
Type & Rank & Parity condition\\\midrule
E1 & 1 & Changes\\
M1 & 1 & Unchanged\\
E2 & 2 & Unchanged\\\bottomrule
\end{tabular}
\[
|J-J'|\le k\le J+J',\qquad m'=m+q,\quad |q|\le k.
\]
E1/M1: \(\Delta J=0,\pm1\), excluding \(0\leftrightarrow0\).
E2: \(\Delta J=0,\pm1,\pm2\) \emph{subject to the triangle inequality};
in particular \(0\leftrightarrow0\) and \(0\leftrightarrow1\) fail.
\subsection*{Orbital screening (one electron)}
E1: \(\Delta l=\pm1\). E2:
\(\Delta l=0,\pm2\), but still apply rank-2
triangle: \(0\leftrightarrow0\) fails.
M1: commonly \(\Delta l=0\) in the leading approximation.
\subsection*{Immediate \(s\)-state examples}
\(s\to p\): possible E1, not E2.

\(s\to d\): possible E2, not E1.
\begin{methodbox}
\textbf{Method:} write operator rank and parity;
check initial/final parity; triangle on \(J\)
(or orbital \(l\) if applicable); projection \(q\);
then evaluate matrix elements.
\end{methodbox}
\end{multicols}
